% Drop-in section for the Tessaris Pilot / COO Mission Layer technical handover.
% Compatible with the handover preamble dated 19 September 2026.
% Uses hyperref (for \path), enumitem and longtable from that document.

\section{Native Pilot Operating Team v2}
\label{sec:pilot-operating-team-v2}

\subsection{Scope and architectural decision}

Pilot Operating Team v2 implements the persistent-coworker pattern natively in
AION.  It does not embed, call, vendor, or take a runtime dependency on
OpenBot.  Grok Bot and OpenBot were treated only as comparative product and
architecture references.  AG-UI compatibility remains a deferred adapter
option; it is not part of the canonical internal runtime.

The governing separation is:
\begin{itemize}[leftmargin=*]
\item the released model selected in the Vault interprets requests, reasons,
plans, and summarises;
\item the COO coordinates bounded missions and Department Pilot assignments;
\item Department Pilots own work within their registered departmental
capabilities;
\item AION owns identity, durable state, capability admission, computer access,
approval, execution boundaries, evidence, receipts, and audit.
\end{itemize}
No provider or model name is embedded in the operating-team service.  The
workspace contract reports \texttt{model\_selection=vault\_selected}; therefore
Qwen, GPT-OSS, a connected hosted model, or another accepted model can be
selected without changing the operating-team contracts.

\subsection{Recovered Train Agent foundation}

The earlier Train Agent system was located rather than replaced.  It already
contained an Electron isolated browser, follow-the-cursor/capture-next-click
recording, semantic selector capture, password masking, a browser skill
library, a chain builder, and an n8n-style visual workflow architect.  Its main
limitation was that compiled browser chains ended in dry-run notifications;
there was no governed semantic replay path.

Version 2 retains those surfaces and adds a durable operating-team service,
semantic demonstration contracts, validation, routines, computer capsules,
mission control, completion packs, scoped memory, worker registration,
templates, and a hash-chained audit.  The existing visual workflow architect
remains available beneath the new Train Agent panel.

\subsection{Runtime components}

\begin{longtable}{p{.24\linewidth}p{.70\linewidth}}
\textbf{Component} & \textbf{Responsibility}\\ \hline
Operating-team service & Durable domain service in
\path{pilot_operating_team_service.py} for capsules, demonstrated skills,
routines, missions, completion packs, memory, workers, templates, preflight
decisions, and audit verification.\\
FastAPI router & API surface in \path{pilot_operating_team_api.py}, rooted at
\path{/api/aion/business/pilot-team}.  Converts domain errors into explicit
HTTP 403, 404, or 422 responses.\\
Train Agent UI & Train Agent and Department Pilot surfaces in
\path{pilot_operating_team_ui.js} for setup, viewing, takeover,
demonstrations, routines, mission controls, workers, templates, and
completion-pack summaries.\\
Desktop styles & Responsive presentation in
\path{pilot_operating_team.css} for the Train Agent workbench and compact
Department Computer bar.\\
Electron bridge & Trusted desktop bridge in \path{electron/main.js} and
\path{preload.js} for opening the isolated browser, recorder control, event
forwarding, profile selection, and validation-only semantic replay.\\
Browser wrapper & \path{electron/isolated-browser.html} binds a durable
Electron partition before navigation and resolves semantic targets inside the
webview.\\
\end{longtable}

The router is registered in both \texttt{backend/desktop\_app.py} and
\texttt{backend/main.py}.  The desktop scheduler invokes both the earlier
scheduled-work service and \texttt{run\_due\_routines()}, so the new generalized
routines do not replace existing scheduled work.

\subsection{Durable workspace layout}

Operating-team data is stored beneath the business-container root at:
\begin{center}\small
\path{<business-containers>/<workspace>/pilot/operating_team/}
\end{center}
The runtime creates separate directories for capsules, department workspaces,
skills, routines, routine runs, missions, completion packs, memory, worker
nodes, and templates.  Each department workspace contains scoped files,
artifacts, and downloads directories.

JSON records are written via a temporary sibling and atomic replacement.  A
process-wide re-entrant lock serializes writes in this service.  This protects
against partial local writes, but it is not a distributed transaction system;
multiple backend processes must not write the same workspace without a future
shared lock or database transaction layer.

\subsection{Department Computer Capsules}

Bootstrap creates capsules for COO, Finance, People, Sales, Marketing, Support,
Products \& Services, and Operations.  Each capsule records:
\begin{itemize}[leftmargin=*]
\item a stable capsule and department identity;
\item a distinct persistent Electron browser profile;
\item a department workspace path;
\item controller state and controller identity;
\item deny-unregistered-destination network policy;
\item department-grant-only connector policy;
\item capsule- and mission-scoped filesystem policy;
\item hidden model secrets and mandatory exact approval for external writes.
\end{itemize}

Electron maps the profile identifier to a partition named
\texttt{persist:<profile-id>}.  Cookies and authenticated sessions are
therefore separated by department.  The current desktop implementation uses
one visible isolated-browser window and switches that window between durable
department profiles; it does not run eight simultaneous virtual desktops.

The user can watch work, take control, and return control to the Department
Pilot.  A transfer changes \texttt{control\_state}, current and previous
controller identifiers, reason, timestamp, record hash, and audit history.
Passwords, passkeys, one-time codes, and CAPTCHA work remain human operations
and are not copied into model context.

\subsection{Teach by demonstration v2}

Starting a teaching session opens the selected Department Computer and enables
the recovered recorder.  Recorded browser events are compiled into a semantic
skill rather than a literal mouse macro.  Each step contains an operation,
semantic target, runtime input reference, readiness wait, validation rule,
approval flag, and capture time.  The skill also records its intended outcome,
owning department, source systems, required inputs, failure policy, and
approval boundaries.

Typed field contents are not embedded in the stored semantic step.  Password,
token, secret, payment-card, and file fields are masked by the recorder; the
service also recursively redacts known secret keys.  Input steps refer to a
runtime input instead of retaining the demonstrated value.

The compiler normalises all of the historical approval representations:
\texttt{approval\_required}, \texttt{requires\_approval}, an
\texttt{approval\_required} safety classification, and known consequential
operations.  This corrects an inherited mismatch in which the old recorder
emitted \texttt{requires\_approval} while the first v2 compiler read only
\texttt{approval\_required}.  Send, submit, delete, purchase, publish, upload,
and explicit approval-checkpoint steps are consequently preserved as mandatory
approval boundaries.

New skills begin as \texttt{draft/not\_tested}.  Publication requires all four
checks to pass: safe inputs, resolved selectors, verified outputs, and verified
approval stops.  Train Agent performs selector validation in the isolated
browser and asks the founder to confirm the expected result in a safe account.
Revalidation does not silently publish a failed or partially checked skill.

\subsection{Semantic replay and live-execution boundary}

The desktop bridge currently admits semantic replay only in
\texttt{dry\_run} mode.  It validates that selectors parse, targets exist, and
targets are visible.  The Electron main process rejects a renderer request for
\texttt{mode=live} with
\texttt{live\_execution\_requires\_backend\_gateway}.

This restriction is intentional.  A renderer-provided approval identifier or
hash is not proof of authority.  Autonomous live replay must be admitted by the
server-side External Tool Gateway, bound to the exact current payload, checked
for expiry and capability authority, and completed through an executor that
returns evidence and a receipt.  The existing AION queue and approval-card
layers deliberately stop at \texttt{waiting\_live\_executor}; this work does
not introduce a second or less-governed executor.  Human takeover can be used
for live browser work now.  Model-driven live side effects remain fail-closed
until the governed executor connection is implemented and tested.

\subsection{Generalised routines and event triggers}

A routine belongs to one Department Pilot and uses exactly one trigger:
\begin{enumerate}[leftmargin=*]
\item a schedule expressed as an interval in minutes; or
\item a business event containing an event type and optional exact field
matches.
\end{enumerate}
It also records an authorised input source, expected result, approval policy,
stale-data policy, no-data policy, counters, test state, enabled state, and
trigger-instance idempotency.

Routines are created disabled.  Enabling is refused until the founder-reviewed
safe test confirms current inputs, output format, audit coverage, approval
stops, and explicit failure states.  An event delivery is deduplicated using
its trigger-instance identifier.  A due schedule or matching event creates a
bounded COO mission and delegates one typed assignment to the owning
department.  It does not execute an external side effect directly.

The current scheduler supports interval schedules with a minimum effective
interval of five minutes.  Calendar expressions, missed-run backfill,
distributed scheduler election, and timezone-specific wall-clock recurrence
are not yet implemented and must not be inferred from the stored timezone
field.

\subsection{COO mission room and in-flight control}

A mission is an outcome contract containing founder constraints, requested
deliverables, owner, state, revision, assignments, and an append-only timeline.
Delegations identify the department, outcome, delegating identity, dependency
references, and handoff depth.  The service caps handoff depth at four and
parallel assignment fan-out at eight to prevent uncontrolled agent loops.

Supported controls are start, pause, resume, stop, complete, fail, and
redirect.  Redirect appends a new instruction and revision without deleting
completed history.  Train Agent currently exposes create, pause, resume,
redirect, and stop.  Department execution remains governed by registered
capabilities and does not inherit authority merely because the COO delegated
the work.

\subsection{Completion packs, memory, and evidence}

A completion pack is hash-bound to the corresponding mission version.  It
separates verified facts and source authorities, assumptions, completed
actions, approvals still waiting, unresolved questions, and artifacts.  This
prevents a prepared draft or intended action from being presented as an
executed result.

Memory is divided into authoritative facts, approved preferences, reusable
procedures, and temporary mission context.  An authoritative fact requires a
named source authority.  Temporary context requires a mission identifier.
Every record states that the selected model may not mutate it directly.  These
checks are application-level authority checks; callers must still be
authenticated by the surrounding desktop/session boundary.

\subsection{Worker nodes and reusable templates}

An optional user-owned worker record contains a stable node identifier,
availability, explicit capabilities, restricted outbound policy, and an
inability to bypass approvals.  A heartbeat can report only a subset of the
capabilities previously registered by the founder.  Attempts to self-add a
capability are rejected.  The present implementation provides registration,
status, and least-authority heartbeat contracts; it does not yet implement a
remote work lease, remote executor, or automatic delivery queue.

Skill, routine, and Pilot configurations can be saved as sanitized templates.
Templates never copy credentials, browser sessions, conversation history, or
authoritative business memory.  A skill template creates a fresh draft that
must be validated again.  A routine template creates a disabled,
\texttt{not\_tested} routine.  A Pilot template reproduces policy configuration
through normal workspace bootstrap rather than copying a browser profile.

\subsection{Action preflight and audit integrity}

The action preflight records initiator identity and kind, department, action
type, target, redacted payload, registered-capability decision, approval
requirement, and canonical payload hash.  Sending, publishing, purchasing,
payments, deletion, overwrite, permission changes, production changes, and
legal acceptance require exact approval.  An unregistered capability is
denied even when the action would otherwise be read-only.

Audit events are appended as JSON Lines.  If $E_i$ is event $i$, the chain is
formed conceptually as
\[
 h_i = H(E_i, h_{i-1}), \qquad h_0 = \texttt{GENESIS}.
\]
The verifier recalculates each canonical event hash and checks every previous
hash reference.  Writes are serialized before the head is updated.  This is a
tamper-evident local chain, not a cryptographically signed or externally
anchored ledger.  A privileged attacker able to rewrite the complete audit and
head files could reconstruct a new chain; higher-assurance deployments should
periodically anchor heads outside the host.

\subsection{API summary}

The native router provides endpoints for workspace bootstrap and inspection,
capsule control transfer, demonstration creation and validation, routine
creation/testing/enablement, event dispatch, due-run dispatch, mission
creation/delegation/control, completion packs, scoped memory, worker
registration and heartbeat, template save and instantiation, action preflight,
and audit verification.

The router itself does not introduce a new public authentication mechanism.
It is designed for the local desktop backend and must remain loopback-only or
be mounted behind AION's existing signed desktop-session and device-pairing
controls.  It must not be exposed directly to a LAN or public interface based
only on caller-supplied strings such as \texttt{person:founder}.

\subsection{Verification and operational notes}

The focused implementation suite contains 17 passing tests covering capsule
isolation, takeover and return, secret-free demonstration compilation,
approval-boundary preservation, skill publication, routine admission,
event deduplication, scheduled dispatch, mission controls, completion packs,
memory authority, preflight denial, worker least authority, template reuse,
and the desktop contract.  Python compilation, Electron and Train Agent
JavaScript syntax checks, inline isolated-browser parsing, and scoped diff
checks also passed.

The repository's pytest environment can remain active during final garbage
collection after printing the complete pass summary and writing
\path{docs/rfc/theorems_results.md}.  During this work it was interrupted only
after the full result had been emitted.  This shutdown delay is test-harness
behaviour, not a failing Pilot test.

After changing Python sources, restart the desktop backend.  After changing
Electron or desktop JavaScript, rebuild the desktop package.  Preserve the
substantial unrelated working-tree changes; this implementation was not
committed automatically.

\subsection{Non-regression requirements}

\begin{itemize}[leftmargin=*]
\item Do not replace Vault model selection with a fixed Qwen, GPT-OSS, or hosted
provider identifier.
\item Do not expose browser partitions, credentials, tokens, typed secrets, or
authoritative memory to the selected model.
\item Do not interpret a renderer-supplied approval hash as authority.
\item Do not enable an untested routine or publish an incompletely validated
demonstrated skill.
\item Do not allow a worker heartbeat, template, routine, or COO delegation to
grant a capability that the founder did not register.
\item Do not describe queued, drafted, or approval-waiting work as completed.
\item Do not expose the Pilot operating-team API without the existing signed
session or equivalent authenticated local boundary.
\item Do not claim simultaneous department virtual machines, cron scheduling,
remote worker execution, externally anchored audit, or autonomous live browser
execution until those facilities are separately implemented and verified.
\item Do not introduce an OpenBot dependency.  Keep any future AG-UI support as
an optional boundary adapter around AION's canonical contracts.
\end{itemize}
